\chapter*{APPENDIX A}
\addcontentsline{toc}{chapter}{APPENDIX A}

\begin{center}
\textbf{Provenance of Every Reported Number}
\end{center}

Every number in this report is regenerated by a named script from committed artifacts. Table~\ref{tab:provenance} maps each results section to the artifact under \texttt{reports/metrics/} and the script under \texttt{scripts/} that produced it.

\begin{table}[htbp]
	\caption{Artifacts and Scripts Behind Each Results Section}
	\centering
	\footnotesize
	\renewcommand{\arraystretch}{1.15}
	\setlength{\tabcolsep}{3pt}
	\begin{tabular}{|L{0.14\linewidth}|L{0.4\linewidth}|L{0.34\linewidth}|}
		\hline
		\textbf{Section} & \textbf{Artifact} & \textbf{Script} \\
		\hline
		\ref{sec:res-bench} & \texttt{benchmark\_results.csv}, \texttt{ablation/} & \texttt{run\_benchmark\_study.py}, \texttt{run\_ablation\_study.py} \\
		\hline
		\ref{sec:res-field} & \texttt{calce\_field\_soh/} & \texttt{run\_field\_soh\_study.py} \\
		\hline
		\ref{sec:res-partial} & \texttt{calce\_partial\_soh/} & \texttt{run\_partial\_soh\_study.py} \\
		\hline
		\ref{sec:res-r} & \texttt{calce\_resistance/} & \texttt{run\_resistance\_study.py} \\
		\hline
		\ref{sec:res-rul} & \texttt{calce\_rul\_horizon\_early\_ref/}, \texttt{error\_bands.csv} & \texttt{run\_rul\_horizon\_study.py}, \texttt{build\_error\_bands.py} \\
		\hline
		\ref{sec:res-sparse} & \texttt{calce\_field\_rul/} & \texttt{run\_field\_rul\_study.py} \\
		\hline
		\ref{sec:res-suff} & \texttt{calce\_sufficiency/} & \texttt{run\_sufficiency\_study.py} \\
		\hline
		\ref{sec:res-oxford} & \texttt{oxford/} & \texttt{run\_oxford\_study.py}, \texttt{run\_oxford\_drive\_cycle.py} \\
		\hline
		\ref{sec:res-xds} & \texttt{cross\_dataset/} & \texttt{run\_cross\_dataset\_ study.py} \\
		\hline
		\ref{sec:res-system} & \texttt{data/interim/rig\_*.txt} & \texttt{run\_serial\_demo.py}, \texttt{panel\_demo.py} \\
		\hline
		\ref{sec:res-card} & Console output of the report card & \texttt{health\_report.py} \\
		\hline
		Figs.~\ref{fig:track}, \ref{fig:cells}, \ref{fig:rul} & \texttt{reports/figures/results\_*.png} & \texttt{export\_results.py} \\
		\hline
	\end{tabular}
	\label{tab:provenance}
\end{table}

\chapter*{APPENDIX B}
\addcontentsline{toc}{chapter}{APPENDIX B}

\begin{center}
\textbf{Reproducing the System and the Results}
\end{center}

The repository is \url{https://github.com/Navvu-gityhub/Behavior-Aware-BMS}. The commands below run on a fresh clone with no hardware, no database, no API keys and no dataset download.

\begin{small}
\begin{verbatim}
git clone \
    https://github.com/Navvu-gityhub/Behavior-Aware-BMS.git
cd Behavior-Aware-BMS
pip install -r requirements-dev.txt

python main.py                     # end-to-end run
python scripts/run_serial_demo.py  # emulated rig
python scripts/health_report.py --serial \
    data/interim/rig_stage_b_voltage_verified.txt
python -m pytest tests/ -q         # the test suite
docker compose up                  # API, gateway, client
\end{verbatim}
\end{small}

With the CALCE data downloaded, the report card of a real cell as it read at a given cycle is printed by

\begin{small}
\begin{verbatim}
python scripts/health_report.py --calce \
    data/raw/calce/CS2/Type2/CS2_35.zip --upto-cycle 90
\end{verbatim}
\end{small}

The raw CALCE and NASA data are not redistributed, and download locations are given in the repository.
